% ======================================================================
\section{Discussion}\label{sec:discussion}
% ======================================================================
% v3: reader-oriented discussion. No statement goes beyond Sections 3-9; limitations repeat Section sec:openpoints.

\purpose{Before concluding we step back from the individual terms and ask four questions: what in this calculation is new, what the structure of the result teaches, how the result is to be used and what a numerical evaluation would require, and what its limitations are.}

\subsection{What is new}\label{sec:disc_new}

\begin{itemize}
\item \emph{The coefficient $\mf A^{(3)}_2$ of the dressing operator,} in coordinate space, with the longitudinal cut-offs and the operator ordering required by the cross section (Section~\ref{sec:coefficients}). Its ingredients, the $O(g^3)$ one-gluon components of the light-cone wave function, were known~\cite{Lublinsky:2016meo}; the ordering of its three-charge part, $\mf A^{(1)}_2\mf N_2$, and its consequences are established here.
\item \emph{A complete organisation of the $O(g^4)$ single inclusive spectrum at mid rapidity} into fifteen products of coefficients, with their longitudinal supports (Section~\ref{sec:groups}).
\item \emph{The fixed-order verification of rapidity factorisation with operator-valued charges:} the two logarithms multiply exactly the BFKL evolution of the projectile correlator and the JIMWLK evolution of the target operator, the latter for an arbitrary ordered correlator, and the terms with more charges cancel in the symmetrised sense (Section~\ref{sec:rapidity}).
\item \emph{The account of all poles,} with the identification of which are ultraviolet and which collinear, and the resulting statements on the scale of the coupling and on the scheme of the fragmentation function (Section~\ref{sec:dglap_rc}).
\item \emph{The finite terms,} row by row, for ordered charges and averaged, with the conditions under which they are integrable (Section~\ref{sec:finite}, Appendices~\ref{app:finite} and~\ref{app:tails}).
\item \emph{The role of the non-commutativity of the charges in the finite terms:} contact terms of the same order as the quadratic terms, and the split of the cubic terms into a part proportional to the C-odd three-point function of the projectile, odd in $\v k$, and a part proportional to the two-point function, even in $\v k$ (Section~\ref{sec:fin_singlet}).
\end{itemize}

\subsection{What the structure of the result teaches}\label{sec:disc_lessons}

\paragraph{One longitudinal variable organises everything.} The second gluon has a longitudinal momentum $p^+$ between $\Lambda$ and $\vee$. For $p^+\gg k^+$ it belongs to the projectile and gives the BFKL term; for $p^+\ll k^+$ it is indistinguishable from the target and gives the JIMWLK term; for $p^+\sim k^+$ it is either collinear to the measured gluon, which gives the DGLAP term, or a genuinely new particle in the final state or in a loop, which gives the finite terms. The three evolutions of the problem are the three boundaries of one phase space, and Row~II.IV alone exhibits all of them (Section~\ref{sec:ex_II_IV}).

\paragraph{The measured gluon is a dressed gluon.} Counting dressed instead of bare gluons is what makes the spectrum vanish without a target, what brings the final-state coefficients $\mf B_3$ and $\mf C_1$ into the amplitudes, and what cancels the ultraviolet part of the self-energy of the measured gluon against the same loop taken across the shock wave. As a consequence the coupling of the emission vertex is renormalised by the wave-function renormalisation of the produced gluon alone.

\paragraph{Charges are operators, and this is visible at NLO.} At leading order the commutators of the charges drop out after the average over a colour-neutral projectile, and the picture of classical sources is sufficient. At order $g^4$ they do not: without them the four-charge terms would not reduce to the evolution, the coefficients of the logarithms would not be the BFKL and JIMWLK ones, and the contact terms of the finite part, which the dilute limit requires, would be missing.

\paragraph{Finite terms are defined together with a prescription for how they are integrated.} Several finite terms are not separately integrable over the position of the unobserved gluon. The result consists of the terms \emph{and} of the statements of Section~\ref{sec:fin_together} about which of them stand under a common integral and in which order the integrals are done. These statements are consequences of the physics (a real and a virtual soft gluon; the measurement as a Fourier integral), not additional choices.

\subsection{How the result is to be used; what a numerical evaluation needs}\label{sec:disc_use}

\paragraph{Inputs.} The projectile enters through the two-point function $n(\v x,\v y)$ and, for the terms odd in $\v k$, through the symmetrised three-point function $\sigma(\v x,\v y,\v z)$; for a projectile that is not globally colour invariant the ordered forms of Appendix~\ref{app:finite} have to be used with the full ordered correlators. The target enters through averages of traces of adjoint Wilson lines. At leading order these are adjoint dipoles. The JIMWLK term and the finite terms require averages of products of more than two Wilson lines at different transverse positions (three in the symmetrised cubic terms; up to four in some two-charge terms as they are written), which are not determined by the dipole alone; the C-odd traces with $D^h$ require a target ensemble that is not C-even to contribute.

\paragraph{Evolution.} The logarithms $\log(\vee/k^+)$ and $\log(k^+/\Lambda)$ are to be absorbed into the evolution of the projectile correlator and of the target averages up to the rapidity of the measured gluon (Section~\ref{sec:logsummary}). A description of NLO accuracy throughout would require, in addition, the NLO evolution of both; this is outside the scope of this paper. The finite terms refer to cut-offs in the plus momentum; with another definition of the rapidity cut-off they would differ by non-logarithmic terms.

\paragraph{Scales.} The renormalisation and factorisation scales enter through the running-coupling and DGLAP terms only, Eq.~\eqref{eq:dg_muRmuF}; the local form~\eqref{eq:dg_local} is the convenient one in momentum space, and it has no large-distance tail of its own. The spectrum of hadrons is obtained by convolution with $\MS$ fragmentation functions defined as number densities; with fragmentation functions in the convention of~\cite{Chirilli:2012jd} the term $\Delta$ of Eq.~\eqref{eq:dg_Delta} has to be subtracted.

\paragraph{Integrals.} A finite term is an integral over one momentum fraction and up to six transverse positions (for instance Eq.~\eqref{eq:S-III.V-3}), two of which carry the measured momentum. A numerical evaluation has to (i) perform the Fourier integrals at the measured momentum first, or in a way that is equivalent to this; (ii) combine the members of the pairs A--D, and the three rows of the two-gluon final state, before the integral over the position of the unobserved gluon; (iii) treat the principal values of Row~I.III and the distributions of the instantaneous terms as defined in Section~\ref{sec:nuance_pv}; (iv) use a rotation-invariant regularisation at the collinear point of Row~I.V. A momentum-space form of all finite terms, which would make step (i) trivial, exists at present only in the dilute limit. We have not carried out such an evaluation and therefore make no statement about the numerical size, or the sign, of the NLO correction.

\subsection{Limitations}\label{sec:disc_limits}

The result is subject to the following restrictions, which are assumptions of the framework and not approximations made along the way: the soft gluons are treated in the eikonal approximation, i.e.\ at leading power in $k^+/\vee$ and $\Lambda/k^+$; the projectile is dilute, i.e.\ the result is a fixed-order expansion in the charge density of the projectile, and it does not describe the collision of two dense objects; the theory is pure Yang--Mills theory; the gluon is produced at mid rapidity, far from both the projectile and the target. Within this framework the points listed in Section~\ref{sec:openpoints} are open; the most important are the comparison of the finite part of the virtual correction with the one-loop central-emission vertex in the dilute limit, the disagreement in the parity of the cubic terms with~\cite{Li:2021ntm}, the dependence on the prescription for the vanishing energy denominators of Row~I.III, and the absence of a numerical evaluation.

\begin{summarybox}
New: $\mf A^{(3)}_2$ with its operator ordering; the organisation of the $O(g^4)$ spectrum; the fixed-order verification of BFKL and JIMWLK factorisation with non-commuting charges; the pole account with its consequences for the coupling and the fragmentation scheme; the finite terms with their integrability conditions. To use the result one needs $n$ and $\sigma$ of the projectile, multi-Wilson-line averages of the target, the evolution of both to the rapidity of the measured gluon, and $\MS$ number-density fragmentation functions. No numerical evaluation has been made.
\end{summarybox}
